\documentclass[11pt]{article}
\usepackage[margin=1in]{geometry}
\usepackage{amsmath,amssymb,amsthm,hyperref}
\newtheorem{proposition}{Proposition}
\newtheorem{theorem}{Theorem}
\newtheorem{corollary}{Corollary}
\newcommand{\KL}{\operatorname{KL}}
\newcommand{\supp}{\operatorname{supp}}
\title{Inexact KL-UOT Refinement Without Objective Acceptance Tests}
\author{Working theory note}
\date{October 19, 2026}
\begin{document}
\maketitle

\begin{abstract}
We derive a refinement guarantee that does not require comparing consecutive primal objective values or accepting a prescribed one-edge trial objective.
The proof tracks restricted optimal values, which need not coincide with the objectives of the computed iterates.
A computable restricted suboptimality bound follows from the marginal-response and complementary stationarity residuals already used by PDUOT, together with a restricted dual-feasibility repair.
The refinement theorem permits inexact solves and controlled deletion of positive flows.
The results are exact-arithmetic statements for KL marginal penalties, not a convergence theorem for the unchanged production solver.
Their novelty relative to existing optimization literature has not been established.
\end{abstract}

\section{Problem and error quantities}
Let $a_i,b_j,\rho_a,\rho_b>0$ and $0\leq c_{ij}\leq C_{\max}<\infty$.
Write $r=X\mathbf1$, $q=X^\top\mathbf1$ and
$$\Phi(X)=\langle c,X\rangle+\rho_a\KL(r\mid a)+\rho_b\KL(q\mid b),\qquad X\geq0.$$
For a node-covering support $\mathcal A$, define
$$v(\mathcal A)=\min_{X\geq0,\ \supp(X)\subseteq\mathcal A}\Phi(X).$$
The zero coupling is feasible and the KL penalties are coercive in transported mass, so the minimum exists.
Let $A_0=\sum_i a_i$, $B_0=\sum_j b_j$, and $B_{\mathrm z}=\Phi(0)=\rho_aA_0+\rho_bB_0$.
Then $0\leq v(\mathcal A)\leq B_{\mathrm z}$ for every support.

For a feasible computed $X$ and arbitrary finite potentials $(f,g)$, let
$$\widehat r_i=a_i e^{-f_i/\rho_a},\qquad \widehat q_j=b_j e^{-g_j/\rho_b},\qquad d_{ij}=c_{ij}-f_i-g_j.$$
Let $s\geq0$ satisfy $s_{ij}X_{ij}=0$ on $\mathcal A$ and set
$$w=d-s,\qquad \omega=\|w\|_\infty,\qquad \nu=\max_{(i,j)\in\mathcal A}(-d_{ij})_+.$$
These are the absolute complementary stationarity quantities used in the PDUOT stationarity mode; in particular $\nu\leq\omega$.
The relative native tolerance must be converted using its actual normalization before applying the statements below.

\section{From PDUOT residuals to restricted suboptimality}
\begin{proposition}[Computable objective-error bound]\label{prop:gap}
Set $f'=f-\rho_a\nu/(\rho_a+\rho_b)$ and $g'=g-\rho_b\nu/(\rho_a+\rho_b)$ and let $D$ denote the KL-UOT dual objective.
The pair $(f',g')$ is feasible on $\mathcal A$, and
$$0\leq\Phi(X)-v(\mathcal A)\leq\Phi(X)-D(f',g')\leq E,$$
where
$$E=\rho_a\KL(r\mid\widehat r)+\rho_b\KL(q\mid\widehat q)+|X|\omega+J(\nu),$$
$$J(\nu)=(\rho_a|\widehat r|+\rho_b|\widehat q|)\big(e^{\nu/(\rho_a+\rho_b)}-1\big).$$
Here $|X|=\sum_{ij}X_{ij}$ and $|u|=\sum_i u_i$ for nonnegative vectors.
\end{proposition}
\begin{proof}
The two downward shifts sum to $\nu$ and remove every restricted dual violation.
The Fenchel identity, valid also for infeasible potentials, is
$$\Phi(X)-D(f,g)=\rho_a\KL(r\mid\widehat r)+\rho_b\KL(q\mid\widehat q)+\langle d,X\rangle.$$
Complementarity gives $\langle d,X\rangle=\langle w,X\rangle\leq |X|\omega$.
Direct substitution gives $D(f,g)-D(f',g')=J(\nu)$.
Weak duality on $\mathcal A$ proves the result.
\end{proof}

All terms can be evaluated in $O(|\mathcal A|+m+n)$ operations without a full-space scan.
The nonnegative bound $E$ avoids relying on the signed split-model gap as an objective-error bound.
It is generally conservative; evaluating the repaired feasible-dual gap directly gives a tighter alternative in exact arithmetic.
If $\widehat r_i\geq\ell_a>0$ and $\widehat q_j\geq\ell_b>0$, then
$$E\leq\frac{\rho_a}{\ell_a}\|r-\widehat r\|_2^2+\frac{\rho_b}{\ell_b}\|q-\widehat q\|_2^2+|X|\omega+J(\omega).$$
Indeed, $\KL(u\mid v)\leq\sum_i(u_i-v_i)^2/v_i$ follows from $\log t\leq t-1$.
For bounded responses, this gives $E=O(\|r-\widehat r\|_2^2+\|q-\widehat q\|_2^2+\omega)$ as the residuals vanish.
Small transported marginals can make the bound weak, so a fixed relative native tolerance alone is not a uniform error guarantee across instances.

\paragraph{Primal stationarity and score uncertainty.}
When $r,q>0$, define $f_i^X=-\rho_a\log(r_i/a_i)$, $g_j^X=-\rho_b\log(q_j/b_j)$ and
$$\eta_a=\|f-f^X\|_\infty,\qquad\eta_b=\|g-g^X\|_\infty,\qquad\eta=\eta_a+\eta_b.$$
Then $|\sigma_{ij}-\sigma_{ij}^X|\leq\eta$, where $\sigma=f+g-c$.
For any fixed $\gamma>0$, the restricted projected-gradient residual obeys
$$R_{\mathcal A}^{\gamma}(X):=\gamma^{-1}\|X-[X-\gamma\nabla_{\mathcal A}\Phi(X)]_+\|_\infty\leq\omega+\eta.$$
To see this, use $\nabla_{\mathcal A}\Phi=s+w+(f-f^X)+(g-g^X)$ and $X=[X-\gamma s]_+$, followed by nonexpansiveness of the projection.
This explicitly links the primal stopping measure to the production stationarity residual and log marginal-response mismatch.

\section{Pruning loss without preserving every positive flow}
Let $\overline X=X-Z\geq0$ be the retained coupling after deleting edges, and define $z_a=Z\mathbf1$, $z_b=Z^\top\mathbf1$.
Assume the selected new edge is retained.
For a trial $Y=X+tE_{ij}$, the corresponding retained trial is $\overline Y=\overline X+tE_{ij}$.
\begin{proposition}[Entropy modulus for deletion]\label{prop:prune}
For nonnegative costs,
$$\Phi(\overline Y)-\Phi(Y)\leq L(Z):=\rho_a\sum_i z_{a,i}\big[1+\log(a_i/z_{a,i})\big]_++\rho_b\sum_j z_{b,j}\big[1+\log(b_j/z_{b,j})\big]_+,$$
with the summand interpreted as zero at zero deletion mass.
\end{proposition}
\begin{proof}
For $h_a(u)=u\log(u/a)-u+a$ and $u\geq z>0$, the function $h_a(u-z)-h_a(u)$ decreases with $u$, because its derivative is $\log((u-z)/u)<0$ for $u>z$.
Its maximum is therefore $h_a(0)-h_a(z)=z(1+\log(a/z))$.
Replacing this maximum by its positive part and summing proves the KL bound.
The deleted transport cost is nonpositive, since $c\geq0$.
\end{proof}
Unlike a gradient bound, this estimate remains finite when deletion produces a zero marginal.
It quantifies positive-flow pruning through deleted node masses rather than requiring exact support preservation.
An edgewise threshold alone does not bound total deletion independently of the number and degrees of deleted edges.

\section{Finite refinement from inner accuracy, without objective acceptance}
Suppose the oracle output $X^k$ has positive marginals and satisfies
$$0\leq\Phi(X^k)-v(\mathcal A_k)\leq\delta_k,\qquad R_{\mathcal A_k}^{\gamma}(X^k)\leq\varepsilon_r.$$
If the maximum inactive primal-consistent score is at most $\varepsilon_f$, stop.
Otherwise select an inactive edge $e_k$ with score $\sigma_{e_k}^{X^k}>\varepsilon_f$, and define
$$h_k=\rho_a/r_i^k+\rho_b/q_j^k,\qquad t_k=\sigma_{e_k}^{X^k}/h_k,\qquad G_k=(\sigma_{e_k}^{X^k})^2/(2h_k).$$
The next support may delete old edges but contains $\supp(\overline X^k)\cup\{e_k\}$, where $\overline X^k=X^k-Z^k$.
The trial is a proof device and need not be formed or used as the next warm start.

\begin{theorem}[Inexact refinement with controlled pruning]\label{thm:refine}
Suppose every nonterminal round satisfies
$$\delta_k+L(Z^k)\leq\theta G_k,\qquad 0\leq\theta<1.$$
Then the restricted optimal values satisfy
$$v(\mathcal A_{k+1})\leq v(\mathcal A_k)-(1-\theta)G_k.$$
If $h_k\leq H<\infty$, the number of nonterminal rounds is at most
$$N\leq\left\lfloor\frac{2H\,v(\mathcal A_0)}{(1-\theta)\varepsilon_f^2}\right\rfloor+1,$$
and the output satisfies $R_{\mathcal E}^{\gamma}(X)\leq\max\{\varepsilon_r,\varepsilon_f\}$.
No comparison of $\Phi(X^{k+1})$ with $\Phi(X^k)$ or a trial objective is assumed.
\end{theorem}
\begin{proof}
The one-edge Taylor bound gives $\Phi(Y^k)\leq\Phi(X^k)-G_k$.
The retained trial is feasible on $\mathcal A_{k+1}$, so
$$v(\mathcal A_{k+1})\leq\Phi(\overline Y^k)\leq\Phi(X^k)-G_k+L(Z^k)\leq v(\mathcal A_k)+\delta_k-G_k+L(Z^k).$$
This proves the descent inequality for restricted optimal values.
Each nonterminal round decreases these nonnegative values by more than $(1-\theta)\varepsilon_f^2/(2H)$.
Summing proves the stated conservative finite bound.
At termination, splitting active and inactive coordinates gives $R_{\mathcal E}^{\gamma}(X)=\max\{R_{\mathcal A}^{\gamma}(X),\max_{e\notin\mathcal A}(\sigma_e^X)_+\}$.
\end{proof}

This theorem tolerates nonmonotone computed objectives: the decreasing sequence is $v(\mathcal A_k)$, not necessarily $\Phi(X^k)$.
The current oracle error $\delta_k$ replaces the previous theorem's acceptance condition on the next objective.
It does not eliminate the need to relate inner accuracy to the desired outer accuracy.
Proposition~\ref{prop:gap} supplies the computable choice $\delta_k=E_k$.

\subsection{Summable errors: no per-round decrease requirement}
The preceding relative forcing rule is not necessary.
Let $e_k=\delta_k+L(Z^k)$ and assume only that the accumulated error is finite.
\begin{theorem}[Finite tolerance with summable errors]\label{thm:summable}
Assume every oracle output satisfies the restricted residual bound, has positive marginals, and every nonterminal round selects an inactive edge with primal-consistent score greater than $\varepsilon_f$.
Suppose $h_k\leq H$ and
$$\sum_{k\geq0}\big(\delta_k+L(Z^k)\big)\leq S<\infty.$$
Then refinement reaches full projected-gradient residual at most $\max\{\varepsilon_r,\varepsilon_f\}$ after at most
$$N\leq\left\lfloor\frac{2H\big(v(\mathcal A_0)+S\big)}{\varepsilon_f^2}\right\rfloor+1$$
nonterminal rounds.
Neither a per-round objective acceptance test nor a per-round condition $e_k<G_k$ is required.
\end{theorem}
\begin{proof}
The proof of Theorem~\ref{thm:refine}, before its relative-error assumption is used, gives
$$v(\mathcal A_{k+1})\leq v(\mathcal A_k)+e_k-G_k.$$
Hence, after $N$ nonterminal rounds,
$$\sum_{k=0}^{N-1}G_k\leq v(\mathcal A_0)-v(\mathcal A_N)+\sum_{k=0}^{N-1}e_k\leq v(\mathcal A_0)+S.$$
Every such round has $G_k>\varepsilon_f^2/(2H)$, which proves finite termination and the stated conservative bound.
The final stationarity estimate follows as in Theorem~\ref{thm:refine}.
\end{proof}

An explicit sufficient schedule is $E_k\leq d_0/(k+1)^2$ and $L(Z^k)\leq l_0/(k+1)^2$, giving $S\leq(d_0+l_0)\pi^2/6$.
This schedule is independent of the unknown restricted optimal values and of the one-edge objective trial.
It allows some rounds to make no objective progress, or even increase the restricted optimal value after pruning.
The fixed-metric PDUOT convergence result makes every positive inner error budget attainable on a fixed node-covering support, provided the restricted residual is also driven below its target.
The pruning budget may be realized trivially by deleting no positive flow; compatibility with a prescribed linear support budget is a separate issue.

\begin{corollary}[Asymptotic full-space stationarity]
Consider an infinite run that always selects a maximizer of the positive inactive primal-consistent score, and takes $G_k=0$ when no positive inactive score exists.
Suppose the objective errors and pruning losses are summable, the selected curvatures are uniformly bounded, and $R_{\mathcal A_k}^{\gamma}(X^k)\to0$.
Then $R_{\mathcal E}^{\gamma}(X^k)\to0$.
If every support covers all nodes, boundedness of the objectives and active residuals gives uniformly positive marginals; every accumulation point is a global KL-UOT optimum.
\end{corollary}
\begin{proof}
The telescoping inequality gives $\sum_kG_k<\infty$, hence $G_k\to0$.
The maximal inactive violation obeys $V_{\mathcal A_k}(X^k)\leq\sqrt{2HG_k}$ and therefore vanishes.
Splitting the projected-gradient residual proves the first conclusion.
Coercivity bounds the couplings, and the node-coverage marginal bounds below keep every accumulation point in the differentiability domain.
Continuity and the convex first-order optimality condition prove the final statement.
\end{proof}

\begin{corollary}[Fixed inner-accuracy regime]
Assume $h_k\leq H$ and choose fixed bounds $\delta_k\leq\overline\delta$, $L(Z^k)\leq\overline L$ such that
$$\overline\delta+\overline L\leq\theta\varepsilon_f^2/(2H),\qquad\theta<1.$$
Then the finite-termination conclusion holds without per-round objective acceptance or adaptive trial comparisons.
\end{corollary}
For fixed errors with $\overline\delta+\overline L>0$, this sufficient condition gives an outer tolerance above $\sqrt{2H(\overline\delta+\overline L)/\theta}$.
It explains why a finite inner tolerance cannot automatically support arbitrarily small outer tolerances.

\paragraph{An explicit uniform curvature bound.}
Uniform curvature need not be a separate assumption when the supports cover all nodes, $R_{\mathcal A_k}^{\gamma}(X^k)\leq\varepsilon_r$, and $\delta_k\leq\overline\delta<\infty$.
Indeed, $\Phi(X^k)\leq B:=B_{\mathrm z}+\overline\delta$ even if the computed objectives are not monotone.
Let $M(s,u)\geq s$ solve $M\log(M/s)-M+s=u$ and define
$$U_a=\max_i M(a_i,B/\rho_a)/a_i,\qquad U_b=\max_j M(b_j,B/\rho_b)/b_j.$$
Active stationarity gives $\nabla_e\Phi(X^k)\geq-\varepsilon_r$.
Using a covering edge at each node yields
$$r_i^k\geq a_i\exp\big((-\varepsilon_r-C_{\max}-\rho_b\log U_b)/\rho_a\big),$$
$$q_j^k\geq b_j\exp\big((-\varepsilon_r-C_{\max}-\rho_a\log U_a)/\rho_b\big).$$
Substitution provides a finite data-dependent $H$ valid for every accepted oracle output.
The constants may be exponentially conservative, especially for small penalties or masses.

\paragraph{Using the actual dual-guided scan.}
The theorem also holds if a selected edge has $\sigma_{e_k}-\eta_k>\varepsilon_f$, replacing $G_k$ by $(\sigma_{e_k}-\eta_k)^2/(2h_k)$.
A scan stopping bound $\max_{e\notin\mathcal A}\sigma_e+\eta_k\leq\varepsilon_f$ guarantees the required inactive condition.
When these lower and upper tests leave an ambiguity region, one must tighten the inner solve or scan with primal-consistent potentials.
Small marginal-response errors in an aggregate relative norm do not by themselves remove this ambiguity.

\section{Direct full-problem objective guarantee}
Replacing $\nu$ in Proposition~\ref{prop:gap} by the maximum dual violation over the complete graph produces a feasible full dual and therefore an upper bound on $\Phi(X)-\Phi^*$.
Only the repair term changes; response defects and active complementarity are unchanged.
Thus the same diagnostics can support both an inner objective-error budget and a final full-problem objective-error bound, with different pricing domains.

\section{Implications and remaining implementation gap}
The new results remove the assumption that each computed objective must pass a previous-trial acceptance test and permit positive-flow pruning with an explicit loss budget.
The summable-error theorem additionally removes any per-round requirement to retain a fixed fraction of the one-edge gain.
They remain conditional convergence results; the inner error and pruning budgets must be met.
They do not establish that the default production tolerance automatically meets those budgets.
The production code already has the relevant primal and dual arrays, responses, complementary slack and full-scan reductions, but does not evaluate $E_k$, log-response score uncertainty, or pruning loss budgets for refinement control.
No source-code modifications are made by this note.

For a fixed node-covering support, the proved fixed-metric PDUOT core converges to a restricted KKT point with positive marginals.
Consequently $E_k$ and $\omega+\eta$ can be made arbitrarily small by continuing that core, in exact arithmetic.
This establishes realizability of a positive inner error budget without requiring a primal trial-objective acceptance test.
It does not supply a practical cost bound for the PID-controlled restarted implementation.

Preserving all positive flows sets $L=0$ but can conflict with a hard support budget.
Allowing controlled deletion quantifies that trade-off but does not prove a uniform linear support bound compatible with arbitrary requested accuracy.
A finite-precision theorem also needs certified cost/scan error bounds and conservative evaluation of the scalar error estimates.
The next empirical question is whether native outputs already satisfy useful ratios $(E_k+L_k)/G_k$, or whether these sufficient bounds are too conservative.
\end{document}
